\par\noindent\textbf{(1)} 令 m₂ 向下位移 x 为正，a=ẍ；绳两侧张力记 T₂,T₁。质点方程 m₂g−T₂=m₂a、T₁−m₁g=m₁a。滑轮转矩 (T₂−T₁)R−bω=Iω̇，因无滑动 ω=ẋ/R，从而 (T₂−T₁)R−b ẋ/R=I ẍ/R。
\par\noindent\textbf{(2)} 消去 T₁,T₂ 得 (m₁+m₂+I/R²)ẍ+(b/R²)ẋ=(m₂−m₁)g。令 M\_eff=m₁+m₂+I/R²、γ=b/R²、F=(m₂−m₁)g，初速零时 ẋ(t)=(F/γ)[1−e\textasciicircum{}\{−γ t/M\_eff\}]。阻力使运动趋终端速度 F/γ，而不是恒加速度。
\par\noindent\textbf{(3)} 由质点方程 T₂=m₂(g−a)、T₁=m₁(g+a)，张力差 T₂−T₁=(m₂−m₁)g−(m₁+m₂)a；代入滑轮式核对额外惯量项。总动能 K=M\_eff ẋ²/2，重力势变化率 −F ẋ，满足 d(K−Fx)/dt=−γ ẋ²，轴承耗散恒非负。
\par\noindent\textbf{(4)} 若 b→0，则 (F/γ)(1−e\textasciicircum{}\{−γt/M\_eff\})→Ft/M\_eff，恢复含滑轮惯量的 Atwood 机。若 I→0，惯性降为 m₁+m₂ 但只在 b 保持有限时仍有耗散；若 b 很大，速度快速限制于小 F/γ，不能把稳态极限误当成 t=0 的初始运动。
